The table is ready. That is where the king's invitation begins in today's Gospel. He has prepared the feast for his son. The guests do not have to supply it. They are asked to come.

Yet some turn away. One goes to his farm, another to his business. Ordinary concerns fill their attention until there is no room for the extraordinary gift being offered. Others answer with violence. The parable's judgment is severe. Refusing God is no small matter. But the king's purpose is not defeated: his servants go out again, and the room fills with people gathered from the roads, bad and good together.

No one in that room can boast of having provided the feast. Everyone has been called to receive. And everyone finds other guests there, people whose lives still need to change. God's gift brings them into company.

Now set beside that crowded room a sentence from our second reading. Paul tells the Philippians that he knows hunger. He knows what it is to have enough, and what it is to go without. Here is someone who belongs to Christ, someone who trusts him, and he can still be hungry.

How do these things belong together? God prepares a feast; his servant lacks food. Isaiah promises abundance; Paul speaks about need. What does God's provision mean for someone whose circumstances are still painful?

We can begin by listening carefully to Paul's claim about strength. Christ enables him to endure both plenty and want. Paul is not saying that faith makes every undertaking succeed. He is speaking about remaining faithful when life changes, when there is much to enjoy and when there is very little to live on. Neither condition has to become his master. Christ sustains him in both.

Then comes the turn we must not miss. Immediately after speaking about Christ's strength, Paul praises the Philippians for sharing his distress. Their help was good. Their generosity mattered. Trust in Christ has not made their gift unnecessary.

Saint John Chrysostom notices how carefully Paul holds those truths together. Paul will not let his benefactors imagine that they own him because they have helped him. But he will not let them think their help is worthless either. He draws them close again by thanking them for their fellowship in his suffering. They have entered into his trouble with him.

This changes the way we speak about faith and need. Someone who trusts God may still need groceries, company, or help with a task that has become too heavy. Their courage does not cancel our responsibility. Their patience does not mean that everything is all right. We cannot point to their faith and use it as our reason for walking away.

And if we are the person who needs help, Paul gives us room to receive it. Depending on Christ does not mean becoming a person who depends on nobody. We need not prove our faith by hiding every difficulty. Gratitude for another person's kindness can belong to trust in God. The gift does not make the giver our owner, and receiving it does not make us less worthy of love.

The Gospel's feast now becomes more searching. The guests who return to farm and business are occupied with what they already have. The invitation asks them to receive something they cannot provide for themselves. We too can become so absorbed in keeping our own lives under control that we resist being gathered into a life with God and with others.

Sometimes that resistance appears as refusal to give: another person's trouble would interrupt our plans. Sometimes it appears as refusal to receive: we would rather remain alone than admit we cannot manage. Paul's words open both closed doors. Christ strengthens him, and the Philippians have done well to help. God's provision makes room for fellowship.

Our opening prayer at Mass asks for grace in the good we must do. Even our generosity needs God's help. We may need courage to notice a difficulty we have been avoiding, patience to keep helping after the first enthusiasm fades, or humility to accept care. Asking for grace brings us toward those acts. It does not make them unnecessary.

But a harder question remains. We can give food and still not cure every illness. We can sit beside someone who mourns and still not restore the person who has died. If God's feast is real, why does so much remain unfinished?

Look at the table in the responsorial psalm. It stands amid enemies. The Shepherd accompanies his people through the dark valley. The psalm does not wait until every danger has disappeared before it speaks of God's care. His presence is real within the trouble. The table is there while the journey is still difficult.

Isaiah carries the promise further. At the Lord's feast, death itself is defeated and tears are wiped away. No comfortable afternoon, no full cupboard, no human generosity can accomplish all of that. God's promise is greater than temporary relief. Its fullness lies ahead.

So hope does not ask the mourner to pretend that grief has ended. It does not tell the hungry person that bodily need is beneath our concern. It teaches us to receive the help given now, to give the help we can now, and to await the good that only God can complete. A shared meal matters, even though it cannot abolish death. A faithful visit matters, even though it cannot wipe away every tear.

At the Eucharist, Christ gives more than a reminder to be generous. He gives the nourishment of his own saving life. The one who feeds his Church is the Lord who suffered and rose from the dead. Our hope rests in him. His gift draws us into communion with him and with the people he gathers.

The prayer after Communion concerns the life that this nourishment gives: participation in God's own life. We remain his creatures, yet his gift truly changes us. We are being formed for a communion whose fullness is still to come. Love offered to a neighbor is one present fruit of that life. It is never the price with which we purchase the feast.

This week, let Paul's thanks make one response concrete. Think of a person whose difficulty you know. Ask what would actually help, and offer something you can carry through: food, an errand, or time to listen. Do not make their gratitude a debt. If you are carrying the difficulty, consider telling a trustworthy person what you need. You do not have to turn your suffering into proof that you can live without anyone.

These acts will not finish the world's sorrow. They can make one person's distress a little less lonely. They can give bodily form to the fellowship Paul praises. They can answer an invitation we might otherwise put off because we are busy managing our own affairs.

The table is ready. We come as people who need what Christ gives, and we find other people beside us who need it too. Receive his gift. Let it become care for someone else. The Shepherd is with us in the valley, and his feast leads toward the day when death will no longer take anyone from the table.
